\section{La ley general antes de sus realizaciones por especie}
\label{sec:ley-general-masa-accion-autonoma}
\index{masa!ley operatorial general}
\index{acción!escala absoluta de masa}
\index{torres espectrales!ley global de masa}

La Mecánica Dimensional recibe de HMT dos ramas ya construidas. La primera
proporciona una recta de acción con dos coordenadas orientadas; la segunda
proporciona la genealogía de cada sección persistente. La masa aparece al
hacer actuar el carácter genealógico sobre una sección de la recta de masa.
Esta confluencia es anterior a cualquier caso particular y, por ello, la
escala \(10^{-34}\) pertenece a la ley general y no a una especie aislada.

\subsection{Rama absoluta: del determinante local a la recta de masa}

El lector determinantal local fija
\begin{equation}
 \Sigma_H=\varphi_{\rm HMT}\alpha_{\rm HMT}^{16},
 \qquad \varepsilon_H:=\left\lfloor\log_{10}\Sigma_H\right\rfloor=-34 .
 \label{eq:ley-masa-sigma-menos34-autonoma}
\end{equation}
La potencia dieciséis procede del conúcleo
\(
 \mathbb Z/2\oplus\mathbb Z/2\oplus\mathbb Z/4
\)
del bloque determinantal. Sobre la recta interna de acción
\(\mathcal L_S\) se obtienen las dos secciones
\begin{equation}
 \boldsymbol\hbar_{\rm pre}
 =H_5\,10^{-34}\mathcal U_S,
 \qquad
 \boldsymbol\hbar_{\rm ret}
 =\eta_{\rm ret}\,10^{-34}\mathcal U_S,
 \label{eq:ley-masa-hbar-pre-ret-autonoma}
\end{equation}
y las cartas aditiva y multiplicativa de su desdoblamiento son
\begin{equation}
 \delta_{2w}=H_5-\eta_{\rm ret},
 \qquad
 R_{\rm act}
 =\frac{\boldsymbol\hbar_{\rm pre}}
        {\boldsymbol\hbar_{\rm ret}}
 =\frac{H_5}{\eta_{\rm ret}}>0 .
 \label{eq:ley-masa-razon-accion-autonoma}
\end{equation}

Sea \(t_0\in\mathcal L_T^+\) la duración elemental del calendario TPK y
\(c_{\rm int}\in\mathcal L_C^+\) su velocidad interna. La periodicidad de
108 iteraciones del endomorfismo temporal CT108 fija la frecuencia angular
\(\omega_0=2\pi/(108t_0)\). Las secciones absoluta anterior y posterior al
retorno son, por tanto,
\begin{equation}
 \begin{aligned}
 \mathbf m_\star^{\rm pre}
 &=\boldsymbol\hbar_{\rm pre}\,\omega_0c_{\rm int}^{-2},\\
 \mathbf m_\star^{\rm ret}
 &=\boldsymbol\hbar_{\rm ret}\,\omega_0c_{\rm int}^{-2},
 \end{aligned}
 \qquad
 \frac{\mathbf m_\star^{\rm pre}}
      {\mathbf m_\star^{\rm ret}}=R_{\rm act}.
 \label{eq:secciones-absolutas-masa-autonoma}
\end{equation}
Ambas pertenecen a
\(
 \mathcal L_M=\mathcal L_S\otimes\mathcal L_T^{-1}
 \otimes\mathcal L_C^{-2}
\).
Así se separan con precisión la década, que fija la escala absoluta, y el
cociente del desdoblamiento bidireccional pre-retorno/post-retorno, que
transporta la memoria de la evolución cerrada.

En coordenadas internas positivas
\(t_0=\widehat t_0\mathcal U_T\) y
\(c_{\rm int}=\widehat c\mathcal U_C\), la reducción escalar de cualquier
sección cuya realización física ya esté fijada adopta la forma
\begin{equation}
 \mathbf m_X^\beta
 =10^{-34}\eta_{\rm ret}\,
  \frac{2\pi}{108\widehat t_0}\,\widehat c^{-2}
  a_X^{(0)}R_{\rm act}^{\kappa_X}\,\mathcal U_M .
 \label{eq:ley-escalar-general-masas-autonoma}
\end{equation}
La coordenada \(a_X^{(0)}\) procede del carácter y de la jerarquía armónica;
\(\kappa_X\) procede del lector de ruta. En particular, la década
\(-34\) pertenece a la sección absoluta y no se incorpora a
\(\kappa_X\).

\subsection{Rama genealógica: registro, carácter y jerarquía armónica}

Cada sección superviviente recorre la cadena
\begin{equation}
 \Omega_{\rm HMT}^{\rm surv}
 \longrightarrow\Gamma_{108}^{\rm enr}
 \longrightarrow\mathcal L_{108}
 \longrightarrow\sigma\in\mathbb Z^5
 \longrightarrow(\sigma,\eta)
 \longrightarrow\chi_{\rm full}(\sigma,\eta)\in\mathbb R_{>0}.
 \label{eq:cadena-genealogica-ley-masa-autonoma}
\end{equation}
El registro conserva el orden de los eventos; la firma conserva sus cinco
invariantes enteros; el carácter los evalúa sin borrar su procedencia. Su
prolongación espectral vive en el semigrupo completo
\begin{equation}
 \mathfrak S=\langle90,120\rangle
 =\{90,120\}\cup\{30r:r\ge6\},
 \qquad
 s_h=q_-^h-q_+^h,
 \quad c_h=q_-^h+q_+^h,
 \label{eq:semigrupo-global-masas-autonoma}
\end{equation}
y adopta la forma
\begin{equation}
 \chi_{\rm full}(\sigma,\eta)
 =\chi_0(\sigma)
  \exp\!\left(\sum_{h\in\mathfrak S}\eta_hs_h\right).
 \label{eq:caracter-completo-masas-autonoma}
\end{equation}
Las alturas \(90,120,180,210,240,270,360,420\) son ocho testigos
genealógicos especialmente visibles. La ley pertenece a todo
\(\mathfrak S\); la recurrencia de segundo orden de \((s_h,c_h)\) prolonga
los testigos a la jerarquía armónica completa.

\subsection{Operador universal sobre las \(13\) multisecciones del dominio material}

El atlas aporta \(81\) posiciones, \(56\) familias de ruta, \(13\)
multisecciones y \(324\) rutas orientadas. Para una multisección \(F\), sea
\(\mathcal H_F=\bigoplus_{\Gamma\in F}\mathbb C e_\Gamma\). El carácter
basal y las torres definen el operador positivo
\begin{equation}
 \widehat{\mathcal A}^{(0)}_F e_\Gamma
 =\chi_{\rm full}(\sigma_\Gamma,\eta_\Gamma)e_\Gamma.
 \label{eq:operador-caracter-basal-masas-autonoma}
\end{equation}
Las dos historias conservadas por TPK inducen lectores
\(r\in\{D,\Omega\}\): \(K_D\) actúa sobre el registro cronológico de
108 iteraciones del endomorfismo temporal CT108 y \(K_\Omega\), sobre la palabra
dodecafásica firmada. Si
\(\widehat K_F^r e_\Gamma=\kappa_r(\Gamma)e_\Gamma\), sea
\(F_s\subseteq F\) el conjunto de rutas de la sección o bloque físico \(s\), y
defínanse
\begin{equation}
 \mathcal H_{F,s}:=\bigoplus_{\Gamma\in F_s}\mathbb C e_\Gamma,
 \qquad
 \widehat{\mathcal A}^{(0)}_{F,s}
 :=\left.\widehat{\mathcal A}^{(0)}_F\right|_{\mathcal H_{F,s}},
 \qquad
 \widehat B_{F,s}^r
 :=\left.\widehat K_F^r\right|_{\mathcal H_{F,s}}.
 \label{eq:restricciones-bloque-ley-masas-autonoma}
\end{equation}
La base diagonal de rutas empleada por los certificados vigentes hace invariante
\(\mathcal H_{F,s}\) bajo ambos operadores y convierte
\(\widehat B_{F,s}^r\) en una restricción autoadjunta. Con estos tipos, la ley
general es
\begin{equation}
 \boxed{
 \widehat{\mathbf M}^{\beta,r}_{F,s}
 =\mathbf m_\star^{\rm ret}\otimes
  R_{\rm act}^{\widehat B_{F,s}^r/2}
  \widehat{\mathcal A}^{(0)}_{F,s}
  R_{\rm act}^{\widehat B_{F,s}^r/2},
 \qquad r\in\{D,\Omega\}.}
 \label{eq:ley-operatorial-general-masas-autonoma}
\end{equation}
La escritura simétrica preserva positividad sin presuponer conmutación. En la
base simultáneamente diagonal empleada por los
certificados vigentes se reduce a
\begin{equation}
 \widehat{\mathbf M}^{\beta,r}_F
 =\mathbf m_\star^{\rm ret}\otimes
  \widehat{\mathcal A}^{(0)}_F
  \exp\!\left((\log R_{\rm act})\widehat K_F^r\right).
 \label{eq:ley-operatorial-diagonal-masas-autonoma}
\end{equation}
La ecuación conserva, en un solo objeto, la escala absoluta, el carácter
global, toda la jerarquía de torres y el corrector bidireccional dependiente
de ruta. La carta física de cada sector actúa después sobre el espectro de
\eqref{eq:ley-operatorial-general-masas-autonoma}. El sector nulo posee su
propio morfismo a cero; los registros compuestos se concatenan antes de
evaluar la firma; los sectores topológicos publican primero sus invariantes
de fusión y trenzado.

\begin{theorem}[Covariancia y cancelación de la escala común]
\label{thm:covariancia-ley-general-masas-autonoma}
La ley~\eqref{eq:ley-operatorial-general-masas-autonoma} es positiva y
covariante bajo cambio unitario de base en cada fibra. Para dos autosecciones
\(X,Y\) realizadas en la misma carta de masa,
\begin{equation}
 \frac{m_Y}{m_X}
 =\frac{\chi_{\rm full}(\sigma_Y,\eta_Y)}
        {\chi_{\rm full}(\sigma_X,\eta_X)}
  R_{\rm act}^{\kappa_r(\Gamma_Y)-\kappa_r(\Gamma_X)}.
 \label{eq:cociente-ley-general-masas-autonoma}
\end{equation}
Los factores comunes \(10^{-34}\mathcal U_S\), \(2\pi/(108t_0)\) y
\(c_{\rm int}^{-2}\) fijan la sección absoluta y se cancelan exactamente en
el cociente.
\end{theorem}

\begin{proof}
El operador \(\widehat{\mathcal A}^{(0)}_F\) es positivo porque sus
autovalores pertenecen a \(\mathbb R_{>0}\); la exponencial de un operador
autoadjunto es positiva. La conjugación por una matriz unitaria preserva el
espectro y, por tanto, la ley. Al dividir dos autovalores realizados en la
misma carta, la sección común \(\mathbf m_\star^{\rm ret}\) se simplifica y
queda~\eqref{eq:cociente-ley-general-masas-autonoma}.
\end{proof}

La fibra electrónica central tiene rango uno y constituye el primer
corolario escalar de esta ley. Las restantes multisecciones conservan el par
de operadores, sus espectros y el morfismo físico que selecciona sabor,
color, orientación o composición. Esta jerarquía permite publicar juntos el
resultado algebraico, su realización dimensional y el contraste externo sin
confundir sus dominios.
